\section{The Render Graph}
\label{sec:rendergraph}

\subsection{Why a Graph}

Three requirements make a hand-written sequence of encoder calls untenable.

First, the aging stage injects parameters at four different pipeline points
(Section~\ref{sec:aging}), so the "pipeline" is not a linear sequence of UI
panels. Second, correctly invalidating the baked LUT
(Section~\ref{sec:lutbake}) and the halation pyramid requires knowing which
parameters each pass depends on, and maintaining that knowledge by hand across
25 parameters and 14 passes is a defect generator. Third, the roadmap calls for
node-based editing, which is a graph by definition; building the graph now costs
little and building it later costs a rewrite.

\subsection{Graph Model}

A graph is a set of \emph{passes}, each declaring:

\begin{itemize}
\item \textbf{Reads}: resource identifiers it samples or reads.
\item \textbf{Writes}: resource identifiers it produces, with descriptors.
\item \textbf{Parameter dependencies}: the set of parameter keys whose values
affect its output.
\item \textbf{Execute}: a closure that encodes GPU work given resolved
resources.
\end{itemize}

\begin{lstlisting}[language=Swift, caption={Pass declaration.}, label={lst:pass}]
protocol RenderPass {
    var name: String { get }
    var reads: [ResourceID] { get }
    var writes: [ResourceWrite] { get }
    /// Keys whose mutation dirties this pass's output.
    var parameterKeys: Set<ParameterKey> { get }

    func encode(into: MTLCommandBuffer,
                resources: ResourceTable,
                params: ResolvedParameters) throws
}

struct ResourceWrite {
    let id: ResourceID
    let descriptor: TextureDescriptorSpec   // resolution-relative
    let access: Access                      // .write, .readWrite
}
\end{lstlisting}

Resource descriptors are \emph{resolution-relative}: a pass declares that it
writes a texture at scale $1$, $1/2$, or $1/16$ of the render extent, and the
concrete size is resolved at compile time from the target extent. This is what
allows the identical graph to serve a $1024$-pixel preview and a $4032$-pixel
export with no branching, and it is the structural guarantee behind NFR-07.

\subsection{Compilation}

Graph compilation runs on the render actor and proceeds in five phases:

\begin{enumerate}
\item \textbf{Topological sort.} Kahn's algorithm over the read/write edges.
A cycle is a programming error and traps in debug builds.
\item \textbf{Dead pass elimination.} Passes whose outputs are not
transitively read by the final output are removed. This is how disabling
halation, grain, or aging removes their cost entirely rather than executing them
with a zero strength --- which is a real and common performance bug in filter
chains.
\item \textbf{Fusion.} Adjacent passes that are pointwise, write to a single
resource of identical extent, and have exactly one consumer are merged into a
single tile dispatch. The fusion candidates are marked by a
\texttt{isPointwise} capability on the pass.
\item \textbf{Lifetime analysis.} For each resource, the first-writer and
last-reader indices in the sorted order define a live interval.
\item \textbf{Allocation.} Linear-scan register allocation over the heap using
the live intervals, followed by barrier insertion at each write-after-read and
read-after-write boundary.
\end{enumerate}

Compilation is memoized on a hash of the enabled-pass set and the target extent.
In steady-state editing --- dragging a slider --- no recompilation occurs;
only the parameter buffers change.

\subsection{Dirty Tracking}

The parameter model is a flat dictionary from \texttt{ParameterKey} to value,
with a monotonically increasing version stamp per key. Each pass caches the
version vector it last executed with. A pass is dirty if any key in its
\texttt{parameterKeys} has advanced, or if any resource it reads was rewritten.

The dirty set propagates forward transitively. The practical consequences are
worth stating explicitly because they are the entire reason the application
feels responsive:

\begin{table}[htbp]
\caption{Incremental Cost of Parameter Changes}
\label{tab:dirtycost}
\begin{center}
\renewcommand{\arraystretch}{1.15}
\begin{tabularx}{\columnwidth}{@{}>{\raggedright\arraybackslash}Xll@{}}
\toprule
\textbf{Changed parameter} & \textbf{Re-runs} & \textbf{Cost} \\
\midrule
Printer lights, print density & LUT + pointwise & $0.9\times$ \\
Push/pull, dev. time, temp.   & LUT + pointwise & $0.9\times$ \\
Print stock, negative stock   & LUT + pointwise & $0.9\times$ \\
Grain amount                  & grain combine only & $0.2\times$ \\
Grain size, seed              & grain field + combine & $0.4\times$ \\
Halation intensity, bias      & halation combine & $0.3\times$ \\
Halation threshold, radius    & full pyramid + all downstream & $1.0\times$ \\
Coupler activity              & interlayer + downstream & $0.8\times$ \\
Vignetting, distortion        & full graph & $1.0\times$ \\
Exposure compensation         & full graph & $1.0\times$ \\
\bottomrule
\multicolumn{3}{@{}p{0.95\columnwidth}@{}}{\footnotesize Cost relative to a full graph execution. Grain amount --- the most-dragged control in user testing --- costs one fifth of a full render.}
\end{tabularx}
\end{center}
\end{table}

\subsection{Resolution Decoupling}

Preview and export are the same graph at different extents, but the preview must
additionally respond to zoom. Three preview modes exist:

\begin{itemize}
\item \textbf{Fit.} Render extent is the view's pixel size, typically
$1179 \times 884$. The entire graph runs at that extent. Full graph time is
approximately 4\,ms.
\item \textbf{Zoomed, $< 100\%$.} Render extent is the view size; the source is
sampled from an appropriate mip of the decoded RAW. Grain and halation are
computed at render extent with the correct physical scaling of
\eqref{eq:formatscale}, so they remain physically consistent.
\item \textbf{Zoomed, $\geq 100\%$.} A region-of-interest render. The graph is
compiled for the visible tile plus an apron whose width is the maximum spatial
support of any enabled pass --- dominated by halation, which at
$\ell_R = 95\,\mu$m and 100\% zoom is roughly $3\ell$ or 32 pixels, plus the
pyramid's implicit support. The apron computation is derived from the passes
rather than hardcoded, which is essential because enabling a large-radius
diffusion filter would otherwise produce a visible seam at tile boundaries.
\end{itemize}

The apron requirement is stated formally: for a chain of passes with spatial
supports $\rho_i$ pixels at their own scale $\varsigma_i$, the apron at render
extent is

\begin{equation}
\rho_{\mathrm{apron}} = \sum_i \frac{\rho_i}{\varsigma_i},
\end{equation}

summing rather than maximizing because supports compose. For the default
configuration this is 96 pixels, which is a 20\% area overhead on a
$1179\times884$ tile --- acceptable, and far cheaper than rendering the full
frame at 100\%.

\subsection{Preview Progressive Refinement}

During an active gesture, the graph renders at half the fit extent and upscales,
which cuts cost by $4\times$ and is imperceptible while the image is moving. On
gesture end, a full-extent render is scheduled after a 60\,ms debounce. If a new
gesture begins during the debounce, the refinement is cancelled before
encoding.

Grain is the exception: it is disabled during gesture and enabled on refinement,
because half-resolution grain that is then upscaled is visibly wrong and the
transition is more distracting than its absence. This is a deliberate,
documented deviation from strict WYSIWYG during motion only.

\subsection{Export Path}

Export runs the same graph at full extent, with three differences:

\begin{enumerate}
\item Tiling. If the full-extent live set would exceed the memory budget, the
image is split into tiles with the computed apron and the graph is executed per
tile. The tile size is chosen so that the live set fits in 400\,MB, and tiles
are processed serially with the heap reused.
\item The grain lattice is in film coordinates, so tiled grain is seamless by
construction --- there is no tile-boundary discontinuity to reconcile. This is
a direct payoff of the resolution-independence design in
Section~\ref{sec:grain}.
\item The output transform targets the requested export color space rather than
the display's.
\end{enumerate}

\begin{figure}[htbp]
\centering
\begin{tikzpicture}[
  font=\scriptsize,
  box/.style={draw=emLine, rounded corners=2pt, minimum height=6mm, text width=24mm, align=center, fill=emGray},
  res/.style={draw=emBlue!60, fill=emBlue!8, rounded corners=1pt, minimum height=5mm, text width=17mm, align=center},
  ar/.style={-{Latex[length=1.4mm]}, draw=emLine}
]
\node[box] (r) at (0,0) {Recipe (value type)};
\node[box, below=5mm of r] (c) {Graph compile\\\scriptsize sort / DCE / fuse};
\node[box, below=5mm of c] (a) {Lifetime + alias};
\node[box, below=5mm of a] (d) {Dirty set};
\node[box, below=5mm of d] (e) {Encode};

\node[res, right=8mm of c] (m1) {memoized on\\pass set + extent};
\node[res, right=8mm of d] (m2) {version vectors\\per parameter key};
\node[res, right=8mm of e] (m3) {MTLCommand\\Buffer};

\foreach \x/\y in {r/c, c/a, a/d, d/e} \draw[ar] (\x) -- (\y);
\draw[ar, dashed] (m1) -- (c);
\draw[ar, dashed] (m2) -- (d);
\draw[ar] (e) -- (m3);
\end{tikzpicture}
\caption{Render graph execution flow. Compilation is memoized; steady-state slider dragging touches only the dirty set and the encode step.}
\label{fig:graphflow}
\end{figure}

\subsection{Testability}

The graph is a pure data structure until \texttt{encode} is called, which makes
several important properties unit-testable without a GPU:

\begin{itemize}
\item \textbf{Topology.} Given a recipe, assert the compiled pass list. This
catches accidental reordering, which is the highest-severity class of imaging
regression.
\item \textbf{Dead pass elimination.} Assert that disabling halation removes
five passes.
\item \textbf{Aliasing correctness.} Assert that no two resources with
overlapping live intervals share an allocation offset. This is checked
exhaustively by a property-based test over randomly generated graphs.
\item \textbf{Dirty propagation.} Assert the re-run set for each parameter key
matches Table~\ref{tab:dirtycost}. This table is generated \emph{from} the test,
not written by hand, so it cannot drift.
\end{itemize}
